\section{Modelo Logico}

\subsection{Estrutura das Tabelas}

\subsubsection{Tabela: usuarios}
\begin{lstlisting}[caption=Estrutura da tabela usuarios]
CREATE TABLE usuarios (
    id UUID PRIMARY KEY DEFAULT gen_random_uuid(),
    nome VARCHAR(100) NOT NULL,
    email VARCHAR(255) UNIQUE NOT NULL,
    senha_hash VARCHAR(255) NOT NULL,
    telefone VARCHAR(20),
    data_nascimento DATE,
    cpf VARCHAR(14) UNIQUE,
    ativo BOOLEAN DEFAULT true,
    email_verificado BOOLEAN DEFAULT false,
    data_ultimo_acesso TIMESTAMP,
    created_at TIMESTAMP DEFAULT CURRENT_TIMESTAMP,
    updated_at TIMESTAMP DEFAULT CURRENT_TIMESTAMP
);
\end{lstlisting}

\subsubsection{Tabela: grupos}
\begin{lstlisting}[caption=Estrutura da tabela grupos]
CREATE TABLE grupos (
    id UUID PRIMARY KEY DEFAULT gen_random_uuid(),
    nome VARCHAR(100) NOT NULL,
    descricao TEXT,
    modelo_organizacao VARCHAR(20) NOT NULL DEFAULT 'individual',
    configuracao JSONB,
    ativo BOOLEAN DEFAULT true,
    created_at TIMESTAMP DEFAULT CURRENT_TIMESTAMP,
    updated_at TIMESTAMP DEFAULT CURRENT_TIMESTAMP,
    
    CONSTRAINT chk_grupos_modelo CHECK (
        modelo_organizacao IN ('individual', 'renda_conjunta', 'splitwise', 'fundo_comum')
    )
);
\end{lstlisting}

\subsubsection{Tabela: usuario\_grupos}
\begin{lstlisting}[caption=Estrutura da tabela usuario_grupos (relacionamento N:M)]
CREATE TABLE usuario_grupos (
    id UUID PRIMARY KEY DEFAULT gen_random_uuid(),
    usuario_id UUID NOT NULL REFERENCES usuarios(id) ON DELETE CASCADE,
    grupo_id UUID NOT NULL REFERENCES grupos(id) ON DELETE CASCADE,
    papel VARCHAR(20) NOT NULL DEFAULT 'membro',
    permissoes JSONB,
    ativo BOOLEAN DEFAULT true,
    data_entrada TIMESTAMP DEFAULT CURRENT_TIMESTAMP,
    data_saida TIMESTAMP,
    
    CONSTRAINT chk_usuario_grupos_papel CHECK (
        papel IN ('administrador', 'membro', 'dependente', 'convidado')
    ),
    UNIQUE($usuario\_id$, $grupo\_id$)
);
\end{lstlisting}

\subsubsection{Tabela: contas}
\begin{lstlisting}[caption=Estrutura da tabela contas]
CREATE TABLE contas (
    id UUID PRIMARY KEY DEFAULT gen_random_uuid(),
    nome VARCHAR(100) NOT NULL,
    descricao TEXT,
    tipo VARCHAR(30) NOT NULL,
    subtipo VARCHAR(30),
    moeda VARCHAR(3) DEFAULT 'BRL',
    saldo_inicial DECIMAL(15,2) DEFAULT 0.00,
    saldo_atual DECIMAL(15,2) DEFAULT 0.00,
    limite_credito DECIMAL(15,2),
    banco VARCHAR(100),
    agencia VARCHAR(10),
    numero_conta VARCHAR(20),
    usuario_id UUID REFERENCES usuarios(id) ON DELETE CASCADE,
    grupo_id UUID REFERENCES grupos(id) ON DELETE SET NULL,
    compartilhada BOOLEAN DEFAULT false,
    ativa BOOLEAN DEFAULT true,
    cor VARCHAR(7) DEFAULT '#007bff',
    icone VARCHAR(50),
    created_at TIMESTAMP DEFAULT CURRENT_TIMESTAMP,
    updated_at TIMESTAMP DEFAULT CURRENT_TIMESTAMP,
    
    CONSTRAINT chk_contas_tipo CHECK (
        tipo IN ('corrente', 'poupanca', 'investimento', 'cartao_credito', 
                'cartao_debito', 'dinheiro', 'outros')
    ),
    CONSTRAINT chk_contas_saldo CHECK (saldo_atual >= 0 OR tipo = 'cartao_credito')
);
\end{lstlisting}

\subsubsection{Tabela: categorias}
\begin{lstlisting}[caption=Estrutura da tabela categorias]
CREATE TABLE categorias (
    id UUID PRIMARY KEY DEFAULT gen_random_uuid(),
    nome VARCHAR(100) NOT NULL,
    descricao TEXT,
    tipo VARCHAR(20) NOT NULL,
    cor VARCHAR(7) DEFAULT '#6c757d',
    icone VARCHAR(50),
    categoria_pai_id UUID REFERENCES categorias(id) ON DELETE SET NULL,
    usuario_id UUID REFERENCES usuarios(id) ON DELETE CASCADE,
    grupo_id UUID REFERENCES grupos(id) ON DELETE CASCADE,
    sistema BOOLEAN DEFAULT false,
    ativa BOOLEAN DEFAULT true,
    ordem INTEGER DEFAULT 0,
    created_at TIMESTAMP DEFAULT CURRENT_TIMESTAMP,
    updated_at TIMESTAMP DEFAULT CURRENT_TIMESTAMP,
    
    CONSTRAINT chk_categorias_tipo CHECK (
        tipo IN ('receita', 'despesa', 'transferencia')
    ),
    CONSTRAINT chk_categorias_escopo CHECK (
        (usuario_id IS NOT NULL AND grupo_id IS NULL) OR
        (usuario_id IS NULL AND grupo_id IS NOT NULL) OR
        (usuario_id IS NULL AND grupo_id IS NULL AND sistema = true)
    )
);
\end{lstlisting}

\subsubsection{Tabela: transacoes}
\begin{lstlisting}[caption=Estrutura da tabela transacoes]
CREATE TABLE transacoes (
    id UUID PRIMARY KEY DEFAULT gen_random_uuid(),
    descricao VARCHAR(255) NOT NULL,
    valor DECIMAL(15,2) NOT NULL,
    tipo VARCHAR(20) NOT NULL,
    data_transacao DATE NOT NULL,
    data_vencimento DATE,
    conta_id UUID NOT NULL REFERENCES contas(id) ON DELETE RESTRICT,
    categoria_id UUID NOT NULL REFERENCES categorias(id) ON DELETE RESTRICT,
    usuario_id UUID NOT NULL REFERENCES usuarios(id) ON DELETE CASCADE,
    grupo_id UUID REFERENCES grupos(id) ON DELETE SET NULL,
    conta_destino_id UUID REFERENCES contas(id) ON DELETE SET NULL,
    status VARCHAR(20) DEFAULT 'confirmada',
    observacoes TEXT,
    tags TEXT[],
    anexos JSONB,
    recorrencia JSONB,
    divisao_splitwise JSONB,
    created_at TIMESTAMP DEFAULT CURRENT_TIMESTAMP,
    updated_at TIMESTAMP DEFAULT CURRENT_TIMESTAMP,
    
    CONSTRAINT chk_transacoes_tipo CHECK (
        tipo IN ('receita', 'despesa', 'transferencia')
    ),
    CONSTRAINT chk_transacoes_status CHECK (
        status IN ('pendente', 'confirmada', 'cancelada', 'agendada')
    ),
    CONSTRAINT chk_transacoes_valor CHECK (valor != 0),
    CONSTRAINT chk_transacoes_transferencia CHECK (
        (tipo = 'transferencia' AND conta_destino_id IS NOT NULL) OR
        (tipo != 'transferencia' AND conta_destino_id IS NULL)
    )
);
\end{lstlisting}

\subsection{Indices Principais}

\begin{lstlisting}[caption=Indices para otimizacao de performance]
-- Indices para consultas frequentes
CREATE INDEX idx_usuarios_email ON usuarios(email);
CREATE INDEX idx_usuarios_ativo ON usuarios(ativo);

CREATE INDEX idx_contas_usuario ON contas(usuario_id);
CREATE INDEX idx_contas_grupo ON contas(grupo_id);
CREATE INDEX idx_contas_ativa ON contas(ativa);

CREATE INDEX idx_transacoes_conta ON transacoes(conta_id);
CREATE INDEX idx_transacoes_categoria ON transacoes(categoria_id);
CREATE INDEX idx_transacoes_usuario ON transacoes(usuario_id);
CREATE INDEX idx_transacoes_data ON transacoes(data_transacao);
CREATE INDEX idx_transacoes_status ON transacoes(status);

-- Indices compostos para consultas complexas
CREATE INDEX idx_transacoes_usuario_data ON transacoes(usuario_id, data_transacao DESC);
CREATE INDEX idx_transacoes_conta_data ON transacoes(conta_id, data_transacao DESC);
CREATE INDEX idx_contas_usuario_ativa ON contas(usuario_id, ativa);
\end{lstlisting}

\subsection{Constraints e Validacoes}

\begin{lstlisting}[caption=Constraints adicionais de integridade]
-- Validacao de email
ALTER TABLE usuarios ADD CONSTRAINT chk_usuarios_email 
    CHECK (email ~* '^[A-Za-z0-9._%+-]+@[A-Za-z0-9.-]+\.[A-Za-z]{2,}$');

-- Validacao de CPF (formato)
ALTER TABLE usuarios ADD CONSTRAINT chk_usuarios_cpf 
    CHECK (cpf IS NULL OR cpf ~ '^\d{3}\.\d{3}\.\d{3}-\d{2}$');

-- Validacao de saldo para cartao de credito
ALTER TABLE contas ADD CONSTRAINT chk_contas_limite_credito 
    CHECK (limite_credito IS NULL OR limite_credito >= 0);

-- Validacao de hierarquia de categorias (evitar loop)
CREATE OR REPLACE FUNCTION check_categoria_hierarchy() RETURNS TRIGGER AS $$
BEGIN
    IF NEW.categoria_pai_id IS NOT NULL AND 
       NEW.categoria_pai_id = NEW.id THEN
        RAISE EXCEPTION 'Categoria nao pode ser pai de si mesma';
    END IF;
    RETURN NEW;
END;
$$ LANGUAGE plpgsql;

CREATE TRIGGER trg_check_categoria_hierarchy
    BEFORE INSERT OR UPDATE ON categorias
    FOR EACH ROW EXECUTE FUNCTION check_categoria_hierarchy();
\end{lstlisting}
